\documentclass[12pt,a4paper]{article}
\usepackage[spanish,es-nodecimaldot]{babel}
\usepackage[utf8]{inputenc}
\usepackage[T1]{fontenc}
\usepackage{lmodern}
\usepackage[top=1.0cm,bottom=1.0cm,left=1.3cm,right=1.3cm]{geometry}
\usepackage{amsmath,amssymb}
\usepackage[table]{xcolor}
\usepackage{enumitem}
\usepackage[normalem]{ulem}
\usepackage{tikz}

\setlength{\parindent}{0pt}
\setlength{\parskip}{11pt}
\pagestyle{empty}

\newcommand{\cuerpo}{\fontsize{15.4}{18.8}\selectfont}
\newcommand{\paso}[1]{\tikz[baseline=(n.base)]{\node[draw, dashed, circle, inner sep=1.2pt, font=\footnotesize\bfseries] (n) {#1};}}

\AtBeginDocument{%
  \setlength{\abovedisplayskip}{6pt}%
  \setlength{\belowdisplayskip}{6pt}%
  \setlength{\abovedisplayshortskip}{3pt}%
  \setlength{\belowdisplayshortskip}{3pt}%
  \cuerpo
}

\begin{document}
\shorthandoff{<>}

{\fontsize{18}{22}\selectfont\bfseries Inferencias basadas en min cuadrados\par}

Una vez que calculamos a y b para nuestra muestra, la inferencia basada en min cuadrados nos permetiran saber que tan confiable es ese modelo para toda la poblacion y no para una muestra en particular. La inferencia sirve para pasar de los estimadores muestrales \textbf{a} y \textbf{b} a conclusiones sobre los parametros poblacionales $\boldsymbol{\alpha}$ y $\boldsymbol{\beta}$.

\textbf{Sxy (varianza de xy)}, \textbf{Sxx(varianza de x)} me van a permitir estimar la pendiente b y la ordenada al origen a. Con \textbf{Syy(varianza de y)} voy a poder calcular la varianza residual muestral

\begin{minipage}[t]{0.36\linewidth}
\vspace{0pt}
\paso{1} \ \ Calc.-{-}$>$\textbf{Sxy, Sxx, Syy}

\vspace{10pt}
\paso{3} \ \
\begin{minipage}[t]{0.72\linewidth}
\vspace{-16pt}
\[
b = \frac{S_{xy}}{S_{xx}}
\]
\[
a = \bar{Y} - b\bar{X}
\]
\end{minipage}
\end{minipage}\hfill
\begin{minipage}[t]{0.38\linewidth}
\vspace{0pt}
\paso{2} \ \
\begin{minipage}[t]{0.80\linewidth}
\vspace{-14pt}
\[
\bar{y} = \frac{\sum y_i}{n} \qquad \bar{x} = \frac{\sum x_i}{n}
\]
\end{minipage}
\end{minipage}\hfill
\begin{minipage}[t]{0.23\linewidth}
\vspace{0pt}
{\small
$s_e^2 \longrightarrow$ Varianza residual muestral

\vspace{6pt}
$s_e \longrightarrow$ Error estandar de estimacion

\vspace{6pt}
\textbf{n-{-}$>$}cant. puntos
}
\end{minipage}

\vspace{4pt}
\paso{4} \ \
\begin{minipage}[t]{0.42\linewidth}
\vspace{-16pt}
\[
s_e^2 = \frac{S_{xx}\cdot S_{yy} - S_{xy}^2}{n(n-2)\cdot S_{xx}}
\]
\[
s_e = \sqrt{s_e^2}
\]
\end{minipage}
\begin{minipage}[t]{0.30\linewidth}
\vspace{-4pt}
\textbf{n--2 gl}
\end{minipage}

\vspace{2pt}
Intervalos de conf. $\alpha$ y $\beta$

\begin{center}
\begin{minipage}{0.36\linewidth}
\begin{center}
$\alpha$
\[
t = \frac{a-\alpha}{s_e}\cdot \sqrt{\frac{n\cdot S_{xx}}{S_{xx}+(n\cdot\bar{x})^2}}
\]
\end{center}
\end{minipage}
\begin{minipage}{0.26\linewidth}
\begin{center}
$\beta$
\[
t = \frac{b-\beta}{s_e}\cdot \sqrt{\frac{S_{xx}}{n}}
\]
\end{center}
\end{minipage}
\begin{minipage}{0.34\linewidth}
$\longrightarrow$ Se compara con el t.crit
\end{minipage}
\end{center}

\[
a - t_{\frac{\alpha}{2}}\cdot s_e\cdot\sqrt{\frac{n\cdot S_{xx}}{S_{xx}+(n\cdot\bar{x})^2}}
\;<\; \alpha \;<\;
a + t_{\frac{\alpha}{2}}\cdot s_e\cdot\sqrt{\frac{n\cdot S_{xx}}{S_{xx}+(n\cdot\bar{x})^2}}
\]

\[
b - t_{\frac{\alpha}{2}}\cdot s_e\cdot\sqrt{\frac{n}{S_{xx}}}
\;<\; \beta \;<\;
b + t_{\frac{\alpha}{2}}\cdot s_e\cdot\sqrt{\frac{n}{S_{xx}}}
\qquad
\textcolor{red}{\text{\footnotesize Errores máximos de estimación de a y b}}
\]

Las pruebas de hipotesis con \textbf{$\boldsymbol{\beta}$} son muy importantes, pq indica el cambio promedio de \textbf{Y} cuando \textbf{X} aumenta 1 unidad. Si $\beta$=0: \textbf{y} no depende de \textbf{x (recta horizontral)}, si \ $\beta \neq 0$ \ : \textbf{y} depende de \textbf{x} por lo que existe relacion lineal.

Habiendo comprobado q \ $\beta \neq 0$ \ osea que existe relacion lineal, podemos utilizar dicha recta ajustada para realizar estimaciones. Podemos ver q pasa en un punto especifico X0

\end{document}
